% =====================================================================
% 1.2 What we do and what we find
% All numbers come from generated/numbers.tex through \val{...}.
% =====================================================================
\subsection{WHAT WE DO AND WHAT WE FIND}
\label{subsec:approach}

We compare three specifications of the joint distribution of the two equity factors. Each is a combination of marginal distributions and a copula, which makes the comparisons between them interpretable:
\begin{description}
\item[M1] uses the empirical distribution of historical returns, i.e.\ empirical margins and empirical dependence, with no parametric assumption;
\item[M2] assumes a bivariate Gaussian distribution, i.e.\ Gaussian margins and a Gaussian copula;
\item[M3] keeps Gaussian margins but couples them with a $t$-copula, which allows for joint extreme movements.
\end{description}
M3 and M2 differ only in the copula, so their gap isolates the effect of tail dependence; M1 and M2 differ in both components, so their gap measures the overall cost of the Gaussian assumption. Models M2 and M3 are estimated by maximum likelihood on weekly log-returns, and each model is used to simulate $10{,}000$ scenarios of the portfolio loss. In every model the default thresholds are calibrated to reproduce the same default probabilities, so all three models imply the same expected loss: any difference in VaR and ES is attributable to the factor distribution alone. The default probabilities supplied with the portfolio are applied to the one-week horizon as given; Section~\ref{sec:model} shows that, once the thresholds are calibrated, the horizon matters only through the shape of the factor distribution.

We then examine the results from three further angles. First, we re-estimate all models on rolling windows of $500$ daily observations to assess how stable the tail estimates are through time. Second, we compare our model-based unexpected loss with the Basel IRB capital charge, which relies on a single Gaussian factor and an infinitely granular portfolio, and decompose the gap between the two. Third, we add a physical climate-risk layer: using the hazard--exposure--vulnerability framework of the CLIMADA package \citep{AznarSiguan19}, we let winter storms erode the value of collateral, so that climate risk acts on recovery rates and not on default probabilities \citep[question~(xi)]{DeGiorgiFilkova26}.

Our main findings are as follows. First, replacing the Gaussian copula by a $t$-copula raises the $99.9\%$ VaR by only \val{q8.M3.VaR_vs_M2|pct1}\%, although the data clearly reject the Gaussian copula: tail dependence between two factors with Gaussian margins changes the distribution of their weighted sum very little, and the dilution anticipated above is almost complete. Second, the empirical distribution raises VaR by \val{q8.M1.VaR_vs_M2|pct0}\%, but this number rests on very few observations, since the $99.9\%$ quantile of the loss is essentially the loss in the second-worst week of the sample. Third, through time the Gaussian VaR is nearly constant by construction, while the empirical VaR jumps by up to \val{q9.ratio_M1_M2_max|excess_pct0}\% above it whenever a crash week enters the estimation window. Fourth, the IRB capital is \val{q10.ratio_UL_M2|inv_pct0}\% of the unexpected loss under M2, and the gap is explained by correlations, not by name concentration. Fifth, the winter-storm overlay is negligible for this portfolio. Damage ratios from Swiss winter storms are tiny, and with storms independent of the credit cycle the expected loss rises by \val{q11.dEL_exact|ppm} parts per million and ES by about as much. Even if the worst storm years coincided with the worst credit years, ES would rise by less than $0.02\%$ (Section~\ref{sec:climate}).

Our contribution is threefold.
\begin{itemize}
\item We separate the effect of the factor margins from that of their dependence structure while holding default probabilities fixed, which isolates the role of the copula in the tail of credit losses.
\item We quantify the Monte Carlo uncertainty of the $99.9\%$ estimates, which rest on only ten tail scenarios out of $10{,}000$, with a reference run of two million scenarios, and show that a single run of the prescribed size cannot rank M2 and M3.
\item We translate the model comparison into regulatory and climate terms: we decompose the distance between the IRB formula and the model-based unexpected loss one assumption at a time, and derive the first-order effect of a recovery-based climate overlay on VaR and ES in closed form, which we confirm by simulation and use to show why geographical diversification hardly matters when storms are independent of credit risk.
\end{itemize}
